\documentclass[10pt, letterpaper]{article}

% Packages:
\usepackage[
ignoreheadfoot, % set margins without considering header and footer
top=2 cm, % seperation between body and page edge from the top
bottom=2 cm, % seperation between body and page edge from the bottom
left=2 cm, % seperation between body and page edge from the left
right=2 cm, % seperation between body and page edge from the right
footskip=1.0 cm, % seperation between body and footer
% showframe % for debugging
]{geometry} % for adjusting page geometry
\usepackage{titlesec} % for customizing section titles
\usepackage{tabularx} % for making tables with fixed width columns
\usepackage{array} % tabularx requires this
\usepackage[dvipsnames]{xcolor} % for coloring text
\definecolor{primaryColor}{RGB}{0, 79, 144} % define primary color
\usepackage{enumitem} % for customizing lists
\usepackage{fontawesome5} % for using icons
\usepackage{amsmath} % for math
\usepackage[
pdftitle={Bipin Lekhak's CV},
pdfauthor={Bipin Lekhak},
pdfcreator={LaTeX with RenderCV},
colorlinks=true,
urlcolor=primaryColor
]{hyperref} % for links, metadata and bookmarks
\usepackage[pscoord]{eso-pic} % for floating text on the page
\usepackage{calc} % for calculating lengths
\usepackage{bookmark} % for bookmarks
\usepackage{lastpage} % for getting the total number of pages
\usepackage{changepage} % for one column entries (adjustwidth environment)
\usepackage{paracol} % for two and three column entries
\usepackage{ifthen} % for conditional statements
\usepackage{needspace} % for avoiding page brake right after the section title
\usepackage{iftex} % check if engine is pdflatex, xetex or luatex
\usepackage{etoolbox} % for \patchcmd and \AtBeginEnvironment

% Ensure that generate pdf is machine readable/ATS parsable:
\ifPDFTeX
    \input{glyphtounicode}
    \pdfgentounicode=1
    % \usepackage[T1]{fontenc} % this breaks sb2nov
    \usepackage[utf8]{inputenc}
    \usepackage{lmodern}
\fi



% Some settings:
\AtBeginEnvironment{adjustwidth}{\partopsep0pt} % remove space before adjustwidth environment
\pagestyle{empty} % no header or footer
\setcounter{secnumdepth}{0} % no section numbering
\setlength{\parindent}{0pt} % no indentation
\setlength{\topskip}{0pt} % no top skip
\setlength{\columnsep}{0cm} % set column seperation
\makeatletter
\let\ps@customFooterStyle\ps@plain % Copy the plain style to customFooterStyle
\patchcmd{\ps@customFooterStyle}{\thepage}{
\color{gray}\textit{\small Bipin Lekhak - Page \thepage{} of \pageref*{LastPage}}
}{}{} % replace number by desired string
\makeatother
\pagestyle{customFooterStyle}

\titleformat{\section}{\needspace{4\baselineskip}\bfseries\large\color{primaryColor}}{}{0pt}{}[\vspace{1pt}\titlerule]

\titlespacing{\section}{
% left space:
-1pt
}{
% top space:
0.3 cm
}{
% bottom space:
0.2 cm
} % section title spacing

\renewcommand\labelitemi{$\circ$} % custom bullet points
\newenvironment{highlights}{
\begin{itemize}[
    topsep=0.10 cm,
    parsep=0.10 cm,
    partopsep=0pt,
    itemsep=0pt,
    leftmargin=0.4 cm + 10pt
    ]
    }{
\end{itemize}
} % new environment for highlights

\newenvironment{highlightsforbulletentries}{
\begin{itemize}[
    topsep=0.10 cm,
    parsep=0.10 cm,
    partopsep=0pt,
    itemsep=0pt,
    leftmargin=10pt
    ]
    }{
\end{itemize}
} % new environment for highlights for bullet entries


\newenvironment{onecolentry}{
\begin{adjustwidth}{
    0.2 cm + 0.00001 cm
    }{
    0.2 cm + 0.00001 cm
    }
    }{
\end{adjustwidth}
} % new environment for one column entries

\newenvironment{twocolentry}[2][]{
\onecolentry
\def\secondColumn{#2}
\setcolumnwidth{\fill, 4.5 cm}
\begin{paracol}{2}
    }{
    \switchcolumn \raggedleft \secondColumn
\end{paracol}
\endonecolentry
} % new environment for two column entries

\newenvironment{header}{
\setlength{\topsep}{0pt}\par\kern\topsep\centering\linespread{1.5}
}{
\par\kern\topsep
} % new environment for the header

\newcommand{\placelastupdatedtext}{ % \placetextbox{<horizontal pos>}{<vertical pos>}{<stuff>}
\AddToShipoutPictureFG*{ % Add <stuff> to current page foreground
\put(
\LenToUnit{\paperwidth-2 cm-0.2 cm+0.05cm},
\LenToUnit{\paperheight-1.0 cm}
){\vtop{{\null}\makebox[0pt][c]{\small\color{gray}\textit{Last updated in October 2026}\hspace{\widthof{Last updated in October 2026}}
}}} %
} %
} %

% save the original href command in a new command:
\let\hrefWithoutArrow\href

% new command for external links:
\renewcommand{\href}[2]{\hrefWithoutArrow{#1}{\ifthenelse{\equal{#2}{}}{ }{#2 }\raisebox{.15ex}{\footnotesize \faExternalLink*}}}

% one line in the skills section: bold label, then the list
\newcommand{\skillline}[2]{\textbf{#1:} #2}


\begin{document}
    \newcommand{\AND}{\unskip
    \cleaders\copy\ANDbox\hskip\wd\ANDbox
    \ignorespaces
    }
    \newsavebox\ANDbox
    \sbox\ANDbox{}

    \placelastupdatedtext
    \begin{header}
        \textbf{\fontsize{24 pt}{24 pt}\selectfont Bipin Lekhak}

        \vspace{0.3 cm}

        \normalsize
        \mbox{{\color{black}\footnotesize\faMapMarker*}\hspace*{0.13cm}Rochester, NY} %
        \kern 0.25 cm %
        \AND %
        \kern 0.25 cm %
        \mbox{\hrefWithoutArrow{mailto:bl1282@rit.edu}{\color{black}{\footnotesize\faEnvelope[regular]}\hspace*{0.13cm}bl1282@rit.edu}} %
        \kern 0.25 cm %
        \AND %
        \kern 0.25 cm %
        \mbox{\hrefWithoutArrow{https://bipinlekhak.com.np/}{\color{black}{\footnotesize\faLink}\hspace*{0.13cm}bipinlekhak.com.np}}%

        \mbox{\hrefWithoutArrow{https://scholar.google.com/citations?user=z8T6GewAAAAJ}{\color{black}{\footnotesize\faGraduationCap}\hspace*{0.13cm}Google Scholar}}%
        \kern 0.25 cm%
        \AND%
        \kern 0.25 cm%
        \mbox{\hrefWithoutArrow{https://openreview.net/profile?id=~Bipin_Lekhak2}{\color{black}{\footnotesize\faFile*}\hspace*{0.13cm}OpenReview}}%
        \kern 0.25 cm%
        \AND%
        \kern 0.25 cm%
        \mbox{\hrefWithoutArrow{https://www.linkedin.com/in/bipin-lekhak/}{\color{black}{\footnotesize\faLinkedinIn}\hspace*{0.13cm}bipin-lekhak}}%
        \kern 0.25 cm%
        \AND%
        \kern 0.25 cm%
        \mbox{\hrefWithoutArrow{https://github.com/bipin-lekhak}{\color{black}{\footnotesize\faGithub}\hspace*{0.13cm}bipin-lekhak}} %
    \end{header}

    \vspace{0.3 cm - 0.3 cm}


    \section{Summary}
    \begin{onecolentry}
        PhD candidate working on \textbf{causal machine learning for time series}: latent-dynamics
        models that predict how a system---a patient, a tumour under therapy, a physical system observed
        through video---will respond to interventions. My models separate a system's native dynamics from its
        response to treatment, are identifiable, and adapt to each individual from its own history through
        meta-learning. Before the PhD, four years as a machine learning engineer shipping medical imaging and
        forecasting systems.
    \end{onecolentry}

    \vspace{0.2 cm}

    \begin{onecolentry}
        \textbf{Research interests:} causal inference over time, treatment-effect estimation under time-varying
        confounding, latent dynamical systems and neural ODEs, identifiable representation learning,
        meta-learning, machine learning for healthcare.
    \end{onecolentry}


    \section{Education}
    \begin{twocolentry}{
        \textit{Aug 2022 – Present}

        \textit{Expected Fall 2027}}
        \textbf{Rochester Institute of Technology}, Rochester, NY

        \textit{Ph.D. in Computing and Information Sciences}
    \end{twocolentry}

    \vspace{0.10 cm}
    \begin{onecolentry}
        \begin{highlights}
            \item Advisor: Prof. Linwei Wang, Computational Biomedicine Laboratory
            \item Dissertation proposal (defended July 2026): \textit{Learning Dynamics for Forecasting Future
            Trajectories under External Interventions}
        \end{highlights}
    \end{onecolentry}

    \vspace{0.2 cm}

    \begin{twocolentry}{
        \textit{Nov 2014 – Sep 2018}}
        \textbf{Pulchowk Campus, Institute of Engineering, Tribhuvan University}, Nepal

        \textit{B.E. in Electronics and Communication Engineering}
    \end{twocolentry}


    \section{Publications}

    \begin{onecolentry}
        \textit{Manuscripts under review}
    \end{onecolentry}

    \vspace{0.10 cm}

    \begin{samepage}
        \begin{twocolentry}{
            \textit{2026}
            }
            \textbf{iASIDE: Identifiable, Adaptable, and Separable Interventional Dynamics}

            \vspace{0.05 cm}

            % TODO: confirm the author list before sending this CV out
            \mbox{\textbf{\textit{Bipin Lekhak}}, Linwei Wang}
        \end{twocolentry}
    \end{samepage}

    \vspace{0.3 cm}

    \begin{onecolentry}
        \textit{Peer-reviewed conference papers}
    \end{onecolentry}

    \vspace{0.10 cm}

    \begin{samepage}
        \begin{twocolentry}{
            \textbf{ICLR 2026}
            }
            \textbf{ProReGen: Progressive Residual Generation under Attribute Correlations}

            \vspace{0.05 cm}

            \mbox{Ruby Shrestha, Ajay Gopi, Casey Meisenzahl, \textbf{\textit{Bipin Lekhak}}, Linwei Wang}

            \vspace{0.05 cm}

            International Conference on Learning Representations.
            \href{https://openreview.net/forum?id=2LzYaW032Q}{OpenReview}
        \end{twocolentry}
    \end{samepage}

    \vspace{0.3 cm}

    \begin{samepage}
        \begin{twocolentry}{
            \textbf{CinC 2023}
            }
            \textbf{Novel In-Home Cardiac Monitoring for Heart Failure Patients}

            \vspace{0.05 cm}

            \textbf{\textit{Bipin Lekhak}}, Ryan Missel, Dillon Dzikowicz, Solomiya Rachynska, Wojciech Zareba, Linwei Wang

            \vspace{0.05 cm}

            Computing in Cardiology.
            \href{https://ieeexplore.ieee.org/abstract/document/10363960}{IEEE Xplore}
        \end{twocolentry}
    \end{samepage}


    \section{Research Experience}

    \begin{twocolentry}{
        \textit{Rochester, NY}

        \textit{Aug 2022 – Present}}
        \textbf{Graduate Research Assistant}

        \textit{Computational Biomedicine Laboratory, RIT} (PI: Prof. Linwei Wang)
    \end{twocolentry}

    \vspace{0.10 cm}
    \begin{onecolentry}
        \begin{highlights}
            \item \textbf{Interventional latent dynamics (ASIDE / iASIDE).} Developed a latent neural-ODE framework
            for forecasting trajectories under time-varying treatments. It separates a system's native dynamics
            from its intervention-response dynamics, adapts each to a subject's history through amortized
            (feed-forward) meta-learning, and establishes identifiability of the latent interventional dynamics.
            Improves intervention-response prediction most at long horizons and under high data heterogeneity.
            \item \textbf{Benchmarking framework.} Built the evaluation suite behind the work: 10+ causal
            time-series baselines (Causal Transformer, CRN, RMSN, MSM, G-Net, SCIP, TE-CDE, TESAR-CDE, ACTIN)
            on tumour-growth simulation, MIMIC-III ICU data and high-dimensional pendulum video, with and without
            time-varying confounding. Hydra, PyTorch Lightning and MLflow; SLURM sweeps across GPU clusters;
            regression tests that verified a migration from PyTorch 1.13 to 2.x against the old results.
            \item \textbf{Simulation.} Wrote a simulator of damped, forced pendulums under stochastic interventions,
            rendered as image sequences, as a benchmark for interventional dynamics from high-dimensional
            observations.
            \item \textbf{Treatment-effect estimation over time} (work in progress, with S.\ Poudel). Temporal X- and
            R-learners that build on the separable latent dynamics to identify causal treatment effects.
            \item \textbf{In-home cardiac monitoring.} Analysed multi-sensor data from a toilet-seat cardiac monitor
            in an NIH-funded (R01NR018301) study of heart-failure patients after hospital discharge, with the
            University of Rochester Medical Center. Measured adherence and deployment failures, built data
            dashboards, and predicted 21-day hospitalisation from 7-day windows (AUC 0.74).
        \end{highlights}
    \end{onecolentry}


    \section{Industry Experience}

    \begin{twocolentry}{
        \textit{Kathmandu, Nepal}

        \textit{Sep 2018 – Jun 2022}}
        \textbf{Machine Learning Engineer}

        \textit{Fusemachines}
    \end{twocolentry}

    \vspace{0.10 cm}
    \begin{onecolentry}
        \begin{highlights}
            \item Built deep-learning semantic segmentation for medical images with U-Net, Mask R-CNN and
            Detectron.
            \item Reconstructed 3D anatomical structures from 2D images by gradient-free optimisation over
            statistical shape models.
            \item Forecast product sales from historical data with a multi-stage model, each stage encoding a
            different assumption about the sales pipeline (LSTMs, Prophet, gradient-boosted trees); designed and
            tested Bayesian variants in PyMC.
        \end{highlights}
    \end{onecolentry}


    \section{Teaching}

    \begin{twocolentry}{
        \textit{Fall 2025}}
        \textbf{Teaching Assistant}, GCCIS-820 Quantitative Foundations, RIT

        \textit{PhD core course: linear algebra, optimisation and statistical inference}
    \end{twocolentry}

    \vspace{0.10 cm}
    \begin{onecolentry}
        \begin{highlights}
            \item Supervised and graded team projects and reports; built the tool used to form project teams.
        \end{highlights}
    \end{onecolentry}


    \section{Skills}

    \begin{onecolentry}
        \skillline{Methods}{causal inference over time (potential outcomes, time-varying confounding, IPW and
        marginal structural models, meta-learners), latent dynamics and state-space models, neural ODEs/CDEs,
        meta-learning, variational inference, identifiability, Bayesian modelling, time-series forecasting,
        medical image segmentation}
    \end{onecolentry}

    \vspace{0.10 cm}

    \begin{onecolentry}
        \skillline{ML stack}{Python, PyTorch, PyTorch Lightning, torchdiffeq, Hydra, MLflow, Ray Tune,
        scikit-learn, PyMC, Prophet, pandas, NumPy}
    \end{onecolentry}

    \vspace{0.10 cm}

    \begin{onecolentry}
        \skillline{Infrastructure}{SLURM and HPC clusters, Linux, Podman/containers, self-hosted MLflow, uv,
        pytest, Git}
    \end{onecolentry}

    \vspace{0.10 cm}

    \begin{onecolentry}
        \skillline{Other}{R, MATLAB, SQL, LaTeX, Streamlit, Dash}
    \end{onecolentry}

\end{document}
